We perform confirmatory comparisons between the graph-guided and baseline linear prompting pipelines on the primary metrics (citation coverage, hallucination rate, reviewer-style scores). For each pre-registered primary hypothesis (“graph-guided > baseline”) we report two-sided paired inference with effect sizes and 95% confidence intervals: paired t-tests when parametric assumptions (approximate normality of paired differences and absence of extreme outliers) are plausible, Wilcoxon signed-rank tests when symmetry/normality are violated, and paired binary tests (e.g., McNemar or exact sign test) for per-instance proportions. When standard assumptions are uncertain or diagnostics indicate poor coverage we report bootstrap bias-corrected intervals and permutation-based p-values; we assess normality, symmetry, and variance patterns and switch to nonparametric or bootstrap inference as needed \cite{web_thelwall_2017_confidence_intervals_for_normalised_citation_cou,web_kabaila_2010_the_coverage_probabililty_of_confidence_interval,web_kabaila_2011_effect_of_a_preliminary_test_of_homogeneity_of_s}.

We control family-wise error rate for the pre-specified primary tests (Holm/Bonferroni family-wise procedures applied at alpha=0.05) and use the Benjamini–Hochberg procedure to control false discovery rate for secondary and exploratory comparisons, reporting both raw and adjusted p-values so readers can inspect multiplicity effects. Reported effect sizes include paired Cohen's d for parametric analyses, rank-biserial or Cliff's delta for nonparametric comparisons, and absolute and relative differences for proportions; confidence intervals for proportions use Agresti–Coull or bootstrap methods as appropriate. All uncertainty estimates explicitly report sample size and exact two-sided p-values alongside the chosen effect-size metric.

Planned ablations (each compared to the full pipeline under the paired framework) are: evidence-injection (disable retrieved passages), scheduler (replace dependency-respecting ordering), schema-constraint (relax JSON constraints), retrieval-depth (vary k and source mix), parallelization (parallel vs sequential leaf generation), and audit-embedding (omit provenance pointers). Each ablation analysis follows the same inferential decision tree and reporting conventions as the primary comparisons.

Reporting conventions: tables and figures routinely show n, a central-tendency summary and dispersion (mean±SD for approximately normal measures or median/IQR for skewed distributions), exact two-sided p-values, adjusted p-values, and effect sizes with 95% CIs. We release complete result files, analysis scripts, and seeds for reproducibility, and we discuss power limitations and sensitivity to analytic choices, emphasizing effect sizes and uncertainty intervals over binary significance claims \cite{web_thelwall_2017_confidence_intervals_for_normalised_citation_cou,web_kabaila_2010_the_coverage_probabililty_of_confidence_interval}.